\section*{Hoofdstuk \ref{ch:b2:conjugate-additions} --- Geconjugeerde carbonylverbindingen: Michael en Robinson}

\begin{solution}{exo:b2:conjugate-additions:1}
\begin{itemize}
  \item \ce{CH3MgBr}: vooral 1,2 (hard, onomkeerbaar), 1-methylcyclohex-2-een-1-ol.
  \item \ce{(CH3)2CuLi}: 1,4, 3-methylcyclohexanon.
  \item Thiofenol met base: 1,4 (het thiolaat is zacht), 3-(fenylsulfanyl)cyclohexanon.
  \item \ce{LiAlH4}: 1,2, cyclohex-2-een-1-ol.
\end{itemize}
\end{solution}

\begin{solution}{exo:b2:conjugate-additions:2}
Donoren: pentaan-2,4-dion, nitromethaan, diethylpropaandioaat (zure \ce{C-H}). Acceptoren:
propeennitril, methylpropenoaat, but-3-een-2-on (\ce{C=C} geconjugeerd met een acceptorgroep).
\end{solution}

\begin{solution}{exo:b2:conjugate-additions:3}
\ce{C4H9Br + 2Li -> C4H9Li + LiBr}, daarna \ce{2C4H9Li + CuI -> (C4H9)2CuLi + LiI}, in ether bij lage
temperatuur onder inert gas.
\end{solution}

\begin{solution}{exo:b2:conjugate-additions:4}
Diethyl-(3-oxobutyl)propaandioaat, \ce{(EtOOC)2CHCH2CH2COCH3}. Hydrolyse geeft het gesubstitueerde
propaandizuur, dat bij verhitten \ce{CO2} verliest (\cref{prop:b2:enolates-aldol:decarboxylation}):
5-oxohexaanzuur, \ce{HOOCCH2CH2CH2COCH3}.
\end{solution}

\begin{solution}{exo:b2:conjugate-additions:5}
Lithiumdimethylcupraat in ether bij lage temperatuur, daarna waterige opwerking: \omterm{def:b2:conjugate-additions:addition-modes}{geconjugeerde additie}.
Methylmagnesiumbromide, harder, addeert vooral 1,2 en geeft 1-methylcyclohex-2-een-1-ol als
hoofdproduct.
\end{solution}

\begin{solution}{exo:b2:conjugate-additions:6}
\ce{Et3N} haalt een beetje \ce{EtSH} weg en geeft \ce{EtS-}; het thiolaat addeert aan het $\beta$-koolstofatoom, het
gevormde \omterm{def:b2:enolates-aldol:enolate}{enolaat} neemt het proton op van een ander \ce{EtSH}, wat \ce{EtS-} terugvormt. Product:
4-(ethylsulfanyl)butaan-2-on. Zwavel is groot en polariseerbaar, zijn vrije elektronenpaar ligt hoog in energie: een \omterm{def:b2:frontier-orbitals:hard-soft}{zacht
nucleofiel}, onder \omterm{def:b2:frontier-orbitals:control}{orbitaalcontrole}, dus het valt aan waar de \omterm{def:b2:frontier-orbitals:frontier}{LUMO} het grootst is, en de additie is
omkeerbaar, wat ook het stabielere 1,4-adduct begunstigt.
\end{solution}

\begin{solution}{exo:b2:conjugate-additions:7}
Cyanide valt het positievere carbonylkoolstofatoom sneller aan: het cyaanhydrine is het kinetische product. Zijn
vorming is omkeerbaar; bij hogere temperatuur, en met de tijd, komt cyanide weer vrij en addeert het opnieuw aan het
$\beta$-koolstofatoom, wat het adduct geeft dat de sterkere $\pi$-binding van \ce{C=O} behoudt, het thermodynamische product.
\end{solution}

\begin{solution}{exo:b2:conjugate-additions:8}
Met 2-methylcyclohexaan-1,3-dion: het Wieland--Miescher-keton (figuur van de \omterm{def:b2:conjugate-additions:robinson}{robinsonannulatie}).
Met cyclohexanon: 4,4a,5,6,7,8-hexahydronaftaleen-2(3\textit{H})-on, een bicyclisch enon waarvan de
\ce{C=C} eindigt op een koolstofatoom van de ringverbinding, zonder angulaire methylgroep.
\end{solution}

\begin{solution}{exo:b2:conjugate-additions:9}
\ce{CH3COCH2CH2CH2COCH3}: C2 en C6 zijn de carbonylgroepen, vijf koolstofatomen van C2 tot C6 inclusief geteld,
een 1,5-diketon. Knip C3--C4 door: donor, het \omterm{def:b2:enolates-aldol:enolate}{enolaat} van propanon op C3 (beter: ethyl-3-oxobutanoaat, waarvan de
estergroep later verwijderd wordt door hydrolyse en \omterm{def:b2:enolates-aldol:decarboxylation}{decarboxylering}); acceptor, but-3-een-2-on. Omstandigheden:
ethyl-3-oxobutanoaat, but-3-een-2-on, katalytisch natriumethoxide in ethanol; daarna waterig zuur en
verhitten.
\end{solution}

\begin{solution}{exo:b2:conjugate-additions:10}
De \omterm{def:b2:enolates-aldol:aldol}{aldolreactie} maakt een binding \ce{C-C} maar zet een $\pi$-binding van \ce{C=O} om in een $\sigma$-binding \ce{C-O}; de
balans is klein en twee moleculen worden er één, dus $K$ is bescheiden en de \omterm{def:b2:enolates-aldol:aldol}{retro-aldolreactie} verloopt gemakkelijk. De
\omterm{def:b2:conjugate-additions:michael}{michaeladditie} maakt een binding \ce{C-C} ten koste van een zwakkere $\pi$-binding van \ce{C=C} en behoudt beide
\ce{C=O}: ze is duidelijk gunstig. De omgekeerde reactie zou vereisen dat het ketonenolaat het gestabiliseerde propaandioaatanion uitstoot; het
evenwicht ligt zo ver aan de kant van het adduct dat het adduct onder katalytische base blijft bestaan.
\end{solution}

\begin{solution}{exo:b2:conjugate-additions:11}
\ce{C(COOEt)2(CH2CH2COOCH3)2}. Na de eerste \omterm{def:b2:conjugate-additions:michael}{michaeladditie} heeft het centrale koolstofatoom nog één
waterstofatoom tussen twee esters, even zuur als voordien: de base haalt het weg en er volgt een tweede additie.
\end{solution}

\begin{solution}{exo:b2:conjugate-additions:12}
Via het meer gesubstitueerde \omterm{def:b2:enolates-aldol:enolate}{enolaat} (naar het koolstofatoom met de methylgroep): het octalon met de methylgroep
op het koolstofatoom van de ringverbinding, 4a-methyl-4,4a,5,6,7,8-hexahydronaftaleen-2(3\textit{H})-on. Via
het minder gesubstitueerde \omterm{def:b2:enolates-aldol:enolate}{enolaat} (aan de kant van \ce{CH2}): de methylgroep komt op een ringkoolstofatoom naast de
ringverbinding. Het eerste wordt begunstigd onder evenwichtsomstandigheden (een alkoxide in zijn alcohol, of de enamineroute
met een \omterm{def:b2:amines:class}{secundair amine} en zuur), die het thermodynamische, meer gesubstitueerde \omterm{def:b2:enolates-aldol:enolate}{enolaat} geven
(\cref{prop:b2:enolates-aldol:enolate-control}).
\end{solution}

\begin{solution}{pb:b2:conjugate-additions:1}
\textbf{1.} Een cyclohexaanring met \ce{C=O} op C1 en C3 en een methylgroep op C2; het waterstofatoom op C2
is het zuurste.
\textbf{2.} Wegname ervan geeft een anion waarvan de lading gedelokaliseerd is over C2 en beide zuurstofatomen; bij
cyclohexanon deelt maar één zuurstofatoom erin, en de zuurgraad is te zwak om in water te meten.
\textbf{3.} Lading op C2; \ce{C1=C2} met \ce{O^-} op C1; \ce{C2=C3} met \ce{O^-} op C3.
\textbf{4.} Ja: met hydroxide is $K = 10^{14.0 - 5.3} = 10^{8.7}$ (water als geconjugeerd zuur, via
$K_e$); met ethoxide $10^{15.9-5.3} = 10^{10.6}$. Beide zijn volledig.
% ledger: ke:25C, pka:ethanol
\textbf{5.} \ce{C1(OH)=C2(CH3)}: C2 draagt nog één waterstofatoom, en meer heeft een \omterm{def:b2:enolates-aldol:tautomerism}{enol} niet nodig; de methylgroep
vervangt alleen het tweede.
\textbf{6.} Zacht: de lading is over drie atomen gespreid en het anion is gestabiliseerd. Het valt het
$\beta$-koolstofatoom aan (\cref{prop:b2:conjugate-additions:regio}).
\textbf{7.} C2 van het \omterm{def:b2:enolates-aldol:enolate}{enolaat} bindt aan het \ce{CH2}-uiteinde van but-3-een-2-on; het gevormde \omterm{def:b2:enolates-aldol:enolate}{enolaat} neemt een
proton op: 2-methyl-2-(3-oxobutyl)cyclohexaan-1,3-dion.
\textbf{8.} Elk nieuw \omterm{def:b2:enolates-aldol:enolate}{enolaat} neemt een proton op van een molecuul dion (of van het oplosmiddel) en vormt zo
het donorenolaat terug: de base wordt niet verbruikt.
\textbf{9.} Het \omterm{def:b2:enolates-aldol:enolate}{enolaat} is ambident; met een zacht koolstofelektrofiel, onder \omterm{def:b2:frontier-orbitals:control}{orbitaalcontrole}, reageert het
op koolstof, en het \ce{C-C}-adduct behoudt twee bindingen \ce{C=O}, de sterkste combinatie; een
\ce{O}-adduct zou, als het al gevormd werd, terugreageren.
\textbf{10.} Eén: C2, gebonden aan C1, C3, de methylgroep en de keten.
\textbf{11.} Het carbonylkoolstofatoom van de keten en het ringcarbonyl C1 (of C3): \ce{C(=O)}--\ce{CH2}--\ce{CH2}--C2--C1,
vijf koolstofatomen inclusief geteld.
\textbf{12.} Om al het dion, het waardevollere reagens, te verbruiken; maar een lichte overmaat, omdat but-3-een-2-on
zeer giftig is en polymeriseert.
\textbf{13.} Ring-C4 en -C6 (naast C3 en C1), de \ce{CH2} van de keten naast haar carbonylgroep, en de eindstandige
\ce{CH3} van de keten. C2 heeft er geen.
\textbf{14.} Het \omterm{def:b2:enolates-aldol:enolate}{enolaat} van de eindstandige \ce{CH3} valt een ringcarbonyl aan: de gesloten ring telt dat
koolstofatoom, het carbonylkoolstofatoom van de keten, twee \ce{CH2}, C2 en het aangevallen carbonylkoolstofatoom: zes atomen.
\textbf{15.} Het \omterm{def:b2:enolates-aldol:enolate}{enolaat} van de inwendige \ce{CH2} van de keten zou een gespannen vierring sluiten. Een
ringenolaat (C4 of C6) dat de carbonylgroep van de keten aanvalt, sluit ook een zesring, maar een overbrugde: dat \omterm{def:b2:enolates-aldol:aldol}{aldol}
kan niet dehydrateren (de dubbele binding zou op een bruggenhoofd van een klein bicyclisch systeem komen) en reageert terug,
terwijl het gecondenseerde \omterm{def:b2:enolates-aldol:aldol}{aldol} door dehydratatie weggetrokken wordt.
\textbf{16.} Het \omterm{def:b2:enolates-aldol:aldol}{aldol} draagt \ce{OH} op het vroegere carbonylkoolstofatoom van de ring, nu aan de ringverbinding; verlies
van water met een waterstofatoom van het vroegere \ce{CH3}-koolstofatoom (nu $\alpha$ ten opzichte van het ketenketon) geeft de
\ce{C=C} geconjugeerd met dat keton.
\textbf{17.} De nieuwe dubbele binding is geconjugeerd met de carbonylgroep: een E1cB-eliminatie via het
\omterm{def:b2:enolates-aldol:enolate}{enolaat}, begunstigd door verhitten (\cref{prop:b2:enolates-aldol:aldol-equilibrium}).
\textbf{18.} Een verzadigd keton en een $\alpha,\beta$-onverzadigd keton (een cyclohexenon).
\textbf{19.} Het koolstofatoom van de ringverbinding dat de methylgroep draagt: vier verschillende substituenten.
\textbf{20.} Nee: het dion is achiraal, en zijn twee carbonylgroepen worden even vaak aangevallen (ze zijn elkaars spiegelbeeld
ten opzichte van de keten); met achirale reagentia is het product racemisch. Chirale aminekatalysatoren
begunstigen één carbonylgroep en geven één enantiomeer in overmaat (deel van jaar~3).
\textbf{21.} Dion \ce{C7H10O2}: \qty{126.155}{g/mol}; \ce{C4H6O}: \qty{70.091}{g/mol}; \ce{C11H14O2}:
\qty{178.231}{g/mol}.
\textbf{22.} \ce{C7H10O2 + C4H6O -> C11H14O2 + H2O}: C 11, H 16, O 3 aan elke kant.
\textbf{23.} $10.0/126.155 = \qty{0.07927}{mol} = \qty{79.3}{mmol}$.
\textbf{24.} $0.07927 \times 178.231 \times 0.70 = \boldsymbol{\approx \qty{9.89}{g}}$
Wieland--Miescher-keton.
\end{solution}
